% ch15 — Gravitational Wave Extraction
\chapter{Gravitational Wave Extraction}
\label{ch:wave-extraction}

The ultimate observable of a numerical relativity simulation is the
gravitational waveform radiated to future null infinity. This chapter develops
the Newman-Penrose formalism for extracting waveforms from the evolved
spacetime and decomposing them into spin-weighted spherical harmonics.

\section{Newman-Penrose Formalism and Weyl Scalars}
\label{sec:we-np}
\coderef{wave\_extraction}{rs}

We introduce the null tetrad, define the five Weyl scalars, and explain why
$\Psi_4$ encodes the outgoing gravitational radiation.

\section{Spin-Weighted Spherical Harmonics}
\label{sec:we-swsh}

We define the spin-weighted spherical harmonics ${}_{s}Y_{\ell m}$ and their
orthogonality relations, providing the angular basis for mode decomposition.

\section{Mode Decomposition of \texorpdfstring{$\Psi_4$}{Psi4}}
\label{sec:we-modes}

We project $\Psi_4$ onto the ${}_{-2}Y_{\ell m}$ basis on coordinate spheres of
finite extraction radius and discuss the convergence with $\ell_{\max}$.

\section{Strain Integration}
\label{sec:we-strain}

We recover the gravitational-wave strain $h = h_+ - i h_\times$ by a
fixed-frequency integration of $\Psi_4$ and address the low-frequency
drift problem.

\section{Radiated Energy, Momentum, and Angular Momentum}
\label{sec:we-radiated}
\coderef{WaveExtraction}{hs}

We derive the formulae for radiated energy, linear momentum (recoil), and
angular momentum from the mode amplitudes, connecting simulations to
astrophysical observables.

\section{Gravitational-Wave Memory\starmark{}}
\label{sec:we-memory}

We introduce the displacement and spin memory effects as DC offsets in the
strain and discuss their extraction from numerical data.

\paragraph{Key References.}
The Newman-Penrose formalism follows \cite{Newman1962,Wald1984}; extraction
techniques follow \cite{Reisswig2011}; memory effects draw on
\cite{Christodoulou1991,Favata2010}.
